\section{失败样板：为什么第一版 Kubernetes 架构不合适}

Vercel 的第一版曾把 Kubernetes 当作工作负载运行平台。Kubernetes 本身并非错误；问题在于产品承诺是“每个 Git commit 都成为一个快速、不可变、可分享的部署”。如果直接把 deployment 映射为 pod 或静态资源单元，部署数量会随 commit 增长，而 idle reservation、调度延迟和 scale transition 会破坏单位经济学。

\subsection{技术可行不等于经济可行}

假设一天产生 $N_d$ 个 preview deployments，每个 deployment 的空闲资源成本为 $c_{idle}$，实际执行成本为 $c_{run}$，则粗略日成本为
\[
C_{day}=N_d\cdot c_{idle}+\sum_j c_{run,j}.
\]
当大部分 preview 很少被访问时，第一项会主导成本。系统必须让“未访问部署”接近零运行成本，同时在第一次请求到来时快速恢复，才能支撑 per-commit product model。

\lecturefigure{02-kubernetes-economics.png}{Per-commit 产品承诺会放大静态资源分配的成本与延迟。}{本地字幕 01:12--01:42、28:50--30:10；概念重绘。}

\begin{knowledgebox}{读图：问题不是 pod，而是映射关系}
若每个 immutable deployment 都长期绑定独立 pod，preview 数量和运行资源近似线性增长。更合适的架构会把部署表示为 artifact + metadata，在请求到达时复用共享 runtime、cache 和调度资源。Kubernetes 可以作为底层组件，但不能让产品对象与资源对象一一绑定。
\end{knowledgebox}

\begin{warningbox}{“Kubernetes 太重”不是普遍结论}
课堂评价针对 Vercel 当时的超高部署基数、快速 scale-to-zero 和 preview 经济模型。对于长期运行服务、稳定容量和成熟平台团队，Kubernetes 仍可能是合理选择。系统设计应比较 workload shape、控制面规模和单位成本，而不是复述工具偏好。
\end{warningbox}

\teachervoice{Rauch 把这次失败描述得很直接：第一版实现“bleeding money left and right”。高价值 teacher voice 在于，他没有降低“每个 commit 都部署”的产品标准，而是把不合理的体验目标转成新的 scheduler、placement 和 metadata 研究问题。}

\subsection{本章小结}

架构必须服从产品对象的基数与生命周期。Per-commit deployment 需要 artifact-first、shared-runtime 和 fast activation，而不是把每个逻辑部署永久映射成一个静态资源租约。

